% =============================================================
% chapter_war_recovery.tex — Контекст реформ ПТО:
% воєнний час, повоєнне відновлення, гармонізація з ЄС.
% Status: stub — наповнюється у Phase 1 / T1.5.
% =============================================================

\section{\MakeUppercase{Драйвери реформи ПТО: воєнний час, повоєнне відновлення, гармонізація з ЄС}}
\label{sec:war-recovery}

\subsection{Воєнний контекст}
\label{sec:war-recovery-war}

% TODO: швидка перепідготовка для оборонної промисловості (зварники,
% оператори БПЛА, кібербезпека); потреби евакуйованих закладів; виклики
% дистанційного моніторингу за умов цифрової нерівності.

\subsection{Повоєнне відновлення}
\label{sec:war-recovery-rebuilding}

% TODO: реконструкція інфраструктури ПТО; інтеграція ветеранів та ВПО
% у програми перепідготовки; критичні професії для відновлення.

\subsection{Зовнішнє фінансування та гармонізація з ЄС}
\label{sec:war-recovery-eu}

% TODO: EU4Skills, ETF (European Training Foundation) рекомендації
% 2024--2026, наближення до EQAVET, Copenhagen Process, EQF mapping;
% наслідки для національного моніторингу (тривалентність даних: МОН,
% ETF, donor reports).

\subsection{Зв'язок з бібліометричними кластерами}
\label{sec:war-recovery-clusters}

% TODO: кластер 1 -- EQAVET; кластер 2 -- ветерани/ВПО;
% кластери 3--4 -- перепідготовка для критичних галузей.
